\par\noindent\textbf{(1)} 巨配分函数 \(\Xi=\sum_N e^{\beta\mu N}Z_N(T,V)\)。利用 \(\partial_{\beta\mu}\ln\Xi\) 得 \(\langle N\rangle\)，再求导得 \[\operatorname{Var}(N)=\partial_{\beta\mu}\langle N\rangle=k_BT(\partial\langle N\rangle/\partial\mu)_{T,V}\]。它是巨正则数涨落的确切公式。
\par\noindent\textbf{(2)} 固定温度下Gibbs–Duhem式 \(dp=n\,d\mu\)，所以 \((\partial n/\partial\mu)_T=n(\partial n/\partial p)_T=n^2\kappa_T\)。于是 \(\operatorname{Var}(N)=k_BTVn^2\kappa_T\)，归一形式为 \(\operatorname{Var}(N)/\langle N\rangle=nk_BT\kappa_T\)。
\par\noindent\textbf{(3)} 理想气体 \(p=nk_BT\)，有 \(\kappa_T=1/p=1/(nk_BT)\)，所以 \(\operatorname{Var}(N)=\langle N\rangle\)。这正是理想经典气体巨正则粒子数Poisson分布；若直接使用严格固定 \(N\) 的正则系综，\(\operatorname{Var}(N)=0\)，不可混淆系综。
\par\noindent\textbf{(4)} 靠近连续相变临界点，\(\kappa_T\) 可能显著增长，因此密度涨落放大并体现临界散射；有限体积会限制相关长度，且长程相关破坏简单独立粒子直观。以上等式要求均匀单相且巨正则平衡可定义，非平衡相分离过程不能直接代入。
